% Editorial fragment for the selective-reuse theory appendix.
\paragraph{Conditional episode equivalence to candidate-null generation.}
\label{app:candidate_null_episode_equivalence}
Consider two versions of the same trained policy, with identical parameters,
instruction, initial environment state, and initial retained factual state.
The first version retains the repertoire interface but fixes its effective
gate to zero at every replan. The second uses the candidate-null interface.
Assume matching decoder token layouts and masks, solver schedules, execution
prefixes, control and observation timing, stopping rules, and native
conditioning interfaces at each replan. Any wall-clock dependence of execution
or termination is held fixed by the comparison protocol. Couple all sampling, environment,
observation, and controller randomness between the two versions.
Assume further that every retained state reaching an ungated decoder input
is updated by the same candidate-independent function of factual observations
and executed commands. Candidate-dependent bookkeeping may differ only when
it has no route to an ungated decoder input.

At the first replan, joint withdrawal gives identical Gaussian sources,
recollection-grounded requirements, zero candidate conditioning, and matched
native and visual context. The fixed solver therefore produces identical
action chunks and executed prefixes. Identical executed commands under the
coupled environment and controller randomness yield identical next states
and observations. The assumed factual-state update then yields identical
retained inputs for the following replan. Induction proves pathwise equality
of the complete executed trajectories. Consequently every measurable success
criterion on those trajectories has equal probability under the two versions.

This is conditional equivalence to the candidate-null version of the same
trained policy, not an empirical comparison with a separately trained
retrieval-free policy. It does not assert that a learned gate becomes zero
when an entire skill is absent from a nonempty bank. That detection property
requires a bank-wide skill-absence protocol and is separate from the
zero-gate premise used here.
